%----------------------------------------------------------------
%
%  File    :  thesis-title.tex
%
%  Author  :  Keith Andrews, ISDS, TU Graz, Austria
% 
%  Created :  22 Feb 1996
% 
%  Changed :  22 Jan 2021
% 
%----------------------------------------------------------------


% --- Main Title Page ------------------------------------------------


\vspace*{2cm}


\begin{center}
    \begin{spacing}{1.1}
        \Huge\sffamily\bfseries
        RespVis:\\
        A Browser-Based, D3 Extension Library\\
        for Creating Responsive SVG Charts
    \end{spacing}

    \vspace{3cm}

    % \LARGE \sffamily Draft 0.9

    \vspace{3cm}

    {\LARGE\sffamily
        Peter Oberrauner
    }
\end{center}



% --- English Title Page ------------------------------------------------


% based on:
% https://tu4u.tugraz.at/fileadmin/Studierende_und_Bedienstete/Formulare/Diplomarbeit_Vorlage.pdf



\cleardoublepage


\vspace*{-3cm}

\begin{center}
    \includegraphics[height=1cm]{diagrams/tugraz-logo.pdf}

    \vspace{2cm}

    \begin{spacing}{1.1}
        \huge\sffamily\bfseries
        RespVis:\\
        A Browser-Based, D3 Extension Library\\
        for Creating Responsive SVG Charts
    \end{spacing}

    \vspace{2cm}

    {\Large\sffamily Peter Oberrauner B.Sc.}

    \vspace{2cm}

    {\Large\sffamily\bfseries Master's Thesis}

    \vspace{5mm}

    {\small\sffamily to achieve the university degree of}

    \vspace{5mm}

    {\normalsize\sffamily Master of Science}

    \vspace{5mm}

    {\normalsize\sffamily
        Master's Degree Programme: Software Engineering and Management
    }


    \vspace{1cm}

    {\small\sffamily submitted to}

    \vspace{5mm}

    {\large\sffamily Graz University of Technology}



    \vspace{1cm}

    {\small\sffamily Supervisor}

    \vspace{5mm}

    {\normalsize\sffamily
        Ao.Univ.-Prof.\ Dr.\ Keith Andrews \\
        Institute of Interactive Systems and Data Science (ISDS)
    }


    \vspace{1cm}

    {\normalsize\sffamily Villach, \thisdate}



    \vfill

    {\footnotesize\sffamily \copyright ~ Copyright \thisyear{} by Peter Oberrauner,
        except as otherwise noted.}

    {\footnotesize\sffamily This work is placed under a
        Creative Commons Attribution 4.0 International
        (\href{https://creativecommons.org/licenses/by/4.0/}{CC BY 4.0}) licence.}


\end{center}




% --- German Title Page ------------------------------------------------

% todo: is the german title page truly necessary?

\cleardoublepage

\begin{otherlanguage}{austrian}

    \vspace*{-3cm}

    \begin{center}
        \includegraphics[height=1cm]{diagrams/tugraz-logo.pdf}

        \vspace{2cm}

        \begin{spacing}{1.1}
            \huge\sffamily\bfseries
            RespVis:\\
            Eine Browser-Basierte, D3 Erweiterungsbibliothek\\
            zur Erstellung von Responsiven SVG Diagrammen
        \end{spacing}

        \vspace{2cm}

        {\Large\sffamily Peter Oberrauner B.Sc.}

        \vspace{2cm}

        {\Large\sffamily\bfseries Masterarbeit}

        \vspace{5mm}

        {\small\sffamily für den akademischen Grad}

        \vspace{5mm}

        \TODO{Should this also be Master of Science?}
        {\normalsize\sffamily Diplom-Ingenieur}

        \vspace{5mm}

        {\normalsize\sffamily
            Masterstudium: Software Engineering and Management
        }


        \vspace{1cm}

        {\small\sffamily an der}

        \vspace{5mm}

        {\large\sffamily Technischen Universität Graz}



        \vspace{1cm}

        {\small\sffamily Begutachter}

        \vspace{5mm}

        {\normalsize\sffamily
            Ao.Univ.-Prof.\ Dr.\ Keith Andrews \\
            Institute of Interactive Systems and Data Science (ISDS)
        }


        \vspace{1cm}

        {\normalsize\sffamily Villach, \thisdate}


        \vspace{1cm}

        {\small Diese Arbeit ist in englischer Sprache verfasst.}



        \vfill

        {\footnotesize\sffamily \copyright ~ Copyright \thisyear{} von Peter Oberrauner, sofern
            nicht anders gekennzeichnet.}

        {\footnotesize\sffamily Diese Arbeit steht unter der Creative Commons
            Attribution 4.0 International
            (\href{https://creativecommons.org/licenses/by/4.0/}{CC BY 4.0})
            Lizenz.}

    \end{center}

\end{otherlanguage}






% --- Pledge ----------------------------------------------------

\cleardoublepage

\vspace*{2cm}


% adapted from:
% https://tu4u.tugraz.at/fileadmin/Studierende_und_Bedienstete/Formulare/Diplomarbeit_Vorlage.pdf
% https://tu4u.tugraz.at/fileadmin/Studierende_und_Bedienstete/Forms/Diploma_thesis_template.pdf
% Email vom 2. Sept. 2015

% and

% Beschluss der Curricula-Kommission für Bachelor-,
% Master- und Diplomstudien vom 10.11.2008
% Genehmigung des Senates am 1.12.2008



\subsection*{Statutory Declaration}
\noindent
\textit{
    I declare that I have authored this thesis independently, that I have
    not used other than the declared sources / resources, and that I have
    explicitly indicated all material which has been quoted either
    literally or by content from the sources used. The document uploaded
    to TUGRAZonline is identical to the present thesis.}

\vspace{1cm}


\begin{otherlanguage}{austrian}

    \subsection*{Eidesstattliche Erklärung}
    \noindent
    \textit{
        Ich erkläre an Eides statt, dass ich die vorliegende Arbeit
        selbstständig verfasst, andere als die angegebenen Quellen/Hilfsmittel
        nicht benutzt, und die den benutzten Quellen wörtlich und inhaltlich
        entnommenen Stellen als solche kenntlich gemacht habe. Das in
        TUGRAZonline hochgeladene Dokument ist mit der vorliegenden
        Arbeit identisch.}

    \vspace{2cm}

    \noindent
    \parbox[top]{4cm}{
        \begin{center}
            \underline{\hspace*{4cm}} \\
            Date/Datum
        \end{center}
    }
    %
    \hfill
    %
    \parbox[top]{6cm}{
        \begin{center}
            \underline{\hspace*{6cm}} \\
            Signature/Unterschrift
        \end{center}
    }

\end{otherlanguage}
\TODO{Sign only printed version or add signature to the digital version?}




% --- English Abstract ----------------------------------------------------


\cleardoublepage

\vspace*{2cm}

\begin{center}
    {\Large\sffamily\bfseries Abstract}
\end{center}

RespVis is an open-source, browser-based library for rendering responsive information visualizations as SVG documents whose elements are mostly styled and positioned via CSS.
It is implemented as a NodeJS package in TypeScript and designed as an extension of D3, thus offering the complete flexibility of D3's convenient document manipulation API to its users.
RespVis consists of various packages containing modules to render low-level visualization components like axes, legends, and series of graphical primitives and high-level premade visualizations for common chart types like bar charts, line charts, and point charts.
The modules provided by RespVis can be used to create responsive visualizations by either  composing them from low-level components or by responsively changing the configurations of premade, high-level visualizations.
One of the main contributions of RespVis is its custom layouter, which enables the various components of a chart to be positioned via any browser-supported CSS layout mechanism like Grid and Flexbox.



% --- German Abstract ----------------------------------------------------



\cleardoublepage

\vspace*{2cm}


\begin{otherlanguage}{austrian}

    \begin{center}
        {\Large\sffamily\bfseries Kurzfassung}
    \end{center}

    \TODO{Translate abstract into german}


\end{otherlanguage}



